% Consolidated textbook-ch03 from commit ace330777e37fc85a6170f63723b6255c77c8f30; original fragments preserved below.
% 2026-11-24 八批严格复核第1批：以下历史片段已局部修订；旧SHA仍指精简前历史版本，不能当作当前片段校验值。

% --- BEGIN textbook-ch03-choices.tex SHA 00c9e2dc13ea54fdd62fea9d9f86b25fa34b0666 ---
% 第三批：《大学物理简明教程》第5版，第3章选择题3.1.1--3.1.8
% 直接依据用户扫描PDF的印刷页68--69（PDF物理页同号）。
% 题面采用保留模型条件的研究性转述；质量状态：待独立复核。
\subsection{选择题：刚体转动基础}

\sourceproblem{3.1.1}{68}{68}
\problemtext 两个力同时作用于一个有固定转轴的刚体。分别判断：两力均平行转轴、均垂直转轴、两力合力为零、两力合力矩为零，能否推出合力矩或合力必然为零？
\workings
轴向力沿转轴的力矩分量均为零，这里所称力矩是关于固定转轴的轴向力矩分量，而非关于任意点的完整力矩矢量，故①正确；垂直轴的力所产生的轴向力矩可以大小相等方向相反，故②所称“可能为零”正确。合力为零的力偶仍可产生非零合力矩，故③错误。合力矩为零也允许非零合力（例如一个沿轴作用的力），故④错误。正确选项为\textbf{B}。
\crosscheck 力与力矩分别是$\sum\bm F$与$\sum\bm r_i\times\bm F_i$，不能互推零值。

\sourceproblem{3.1.2}{68}{68}
\problemtext 转动惯量为$J$、有效半径$R$的定滑轮绕轻绳旋转。原先绳下端悬挂重力为$P$的物体；若取下物体改为用大小同为$P$的恒力竖直下拉轻绳，比较滑轮角加速度。
\workings
原悬重物质量$P/g$，绳张力$T<P$，联立
$P-T=(P/g)a$、$TR=J\alpha$、$a=R\alpha$得
\[
\alpha_{\mathrm{hang}}=\frac{PR}{J+(P/g)R^2}.
\]
改为外力直接拉绳时$T=P$，故
\[
\alpha_{\mathrm{pull}}=\frac{PR}{J}>\alpha_{\mathrm{hang}}.
\]
原题问“物体去掉而以同等力直接下拉”时角加速度如何变化，答案为\textbf{变大（C）}。
\crosscheck 仅当$J>0$及绳不滑动时上述比较成立；不能误把重力$P$当成原先绳张力。

\sourceproblem{3.1.3}{68}{68}
\problemtext 对刚体的转动惯量$J$，判断它与质量、质量的空间分布、转轴的位置之间的依赖关系。
\workings
固定轴转动惯量定义
\[
J=\int r_\perp^2\,dm.
\]
其数值同时依赖刚体质量、各质量元距轴的垂直距离及转轴选取；只有三种因素都考虑才正确。原题选\textbf{C}。
\crosscheck 同质量的细环与圆盘绕对称轴分别有$MR^2$与$\frac12MR^2$，说明分布重要。

\sourceproblem{3.1.4}{68--69}{68--69}
\problemtext 初始绕光滑固定轴转动的圆盘，被两颗质量相同、速率相同、方向相反、且沿\emph{同一条偏离圆心的直线}射来的子弹射入并留在盘内。比较射入后的盘角速度。
\workings
取盘心为角动量参考点。两颗子弹的入射动量沿同一直线而方向相反，因而其入射角动量代数和为零。极短碰撞中轴的冲量矩为零，
\[
J_0\omega_0=J_f\omega_f,\qquad J_f>J_0,
\]
故在$\omega_0\ne0$时有
\[
\boxed{|\omega_f|=\frac{J_0}{J_f}|\omega_0|<|\omega_0|}.
\]
应选\textbf{C（减小）}。
\crosscheck 若两颗弹丸不是同时按同一直线反向进入，或者圆盘起初静止，应重新核对几何与入射角动量；这里严格使用原题配图的对向共线条件。

\sourceproblem{3.1.5}{69}{69}
\problemtext 转动惯量为$J$的圆盘由$\omega_0$减速到$\omega_0/2$，问全过程阻力矩所做的功。
\workings
转动动能定理：
\[
W=\frac12J\left(\frac{\omega_0}{2}\right)^2-\frac12J\omega_0^2
=\boxed{-\frac38J\omega_0^2}.
\]
原题选\textbf{D}。
\crosscheck 阻力矩作负功，数值单位为焦耳；不应误写为角动量变化。

\sourceproblem{3.1.6}{69}{69}
\problemtext 刚体绕定轴\emph{变速}转动，距轴$r$的质元的法向加速度$a_n$与切向加速度$a_\tau$可能怎样随时间变化？
\workings
\[
a_n=r\omega^2,\qquad a_\tau=r\alpha=r\dot\omega.
\]
因角速率随时变化，$a_n$通常变化；但若角加速度$\alpha$恒定，$a_\tau$可以不变。因此只有“$a_n$变化，$a_\tau$可以不变”的判断成立，选\textbf{C}。
\crosscheck 此题取$r>0$且“变速”指角速率大小变化；距轴为零的特殊质元两项均零。

\sourceproblem{3.1.7}{69}{69}
\problemtext 均匀砂轮（质量$M$、半径$R$）绕中心轴以$\omega$转动，其动能等于质量为$M$的物体从高度$h$自由落地获得的动能。求$h$。
\workings
均匀圆盘转动惯量$J=\frac12MR^2$，故
\[
\frac12J\omega^2=Mgh
\Longrightarrow \boxed{h=\frac{R^2\omega^2}{4g}}.
\]
原题选\textbf{D}。
\crosscheck 半径的平方与角速度平方之积除以$g$，量纲为长度。

\sourceproblem{3.1.8}{69}{69}
\problemtext 长度同为$L$、质量同为$m$的细线单摆和均匀细杆物理摆均由与竖直线成$\theta$的位置静止释放，比较通过最低位置时的角速率。
\workings
单摆：$mgL(1-\cos\theta)=\frac12mL^2\omega_p^2$，即
\[
\omega_p^2=\frac{2g(1-\cos\theta)}{L}.
\]
细杆关于上端的转动惯量$J=\frac13mL^2$，
质心下降$\frac12L(1-\cos\theta)$，故
\[
\frac12mgL(1-\cos\theta)=\frac16mL^2\omega_r^2
\Longrightarrow \omega_r^2=\frac{3g(1-\cos\theta)}{L}.
\]
对通常的轻绳单摆静止释放，初始张力$T_0=mg\cos\theta$，故须$0<\theta\le\pi/2$才可直接采用绷紧的圆弧模型（$\theta=\pi/2$时初张力为零而随后为正）。此范围内$\omega_p<\omega_r$，选\textbf{C}。若$\theta>\pi/2$，轻绳初始松弛，须另建模；若改为刚性轻杆约束，则能量公式在$0<\theta<\pi$可用，精确倒立$\theta=\pi$静止释放为不稳定平衡，不能断言必经过最低点。
\crosscheck 棒的质心在$L/2$，不能把细杆的质量当成位于末端的质点。
% --- END textbook-ch03-choices.tex ---

% --- BEGIN textbook-ch03-blanks.tex SHA 2ba5c7288df840728dde69fef09349a1a6217219 ---
% 第三批：《大学物理简明教程》第5版，3.2.1--3.2.12。
% 源书印刷页69--70，用户扫描PDF物理页同号；质量状态待独立复核。
\subsection{填空题：定轴转动、角动量和能量}

\sourceproblem{3.2.1}{69}{69}
\problemtext 三个各为$m$的质点位于边长$a$的等边三角形顶点。求过重心垂直三角形的轴、过重心且平行任一边的轴、过重心及一个顶点的面内轴的转动惯量。
\workings
顶点到重心的距离均为$a/\sqrt3$，故法向轴
\[
\boxed{J_0=3m(a/\sqrt3)^2=ma^2}.
\]
由平面薄层的垂直轴定理及等边三角形的$120^\circ$旋转对称性，任何一条面内重心轴均有
\[
\boxed{J_1=J_2=\tfrac12J_0=\tfrac12ma^2}.
\]
\crosscheck $J_0=J_x+J_y$；每一质量点的平方距离非负，且两条正交面内重心轴由于三重对称性具有相同转动惯量。

\sourceproblem{3.2.2}{69}{69}
\problemtext 两个厚度、总质量相等而均匀密度不同的圆盘，$\rho_A>\rho_B$；比较绕各自盘心垂直盘面的轴的转动惯量$J_A,J_B$。
\workings
$M=\rho\pi R^2h$；在$M$与厚度$h$固定下，
$R^2=M/(\rho\pi h)$。均匀圆盘$J=\frac12MR^2$，因此
\[
\frac{J_A}{J_B}=\frac{\rho_B}{\rho_A}<1,
\quad \boxed{J_A<J_B}.
\]
\crosscheck 密度较大、质量相同时半径更小，质量更靠近转轴。

\sourceproblem{3.2.3}{69}{69}
\problemtext 刚体绕定轴转动，$J=3.0\,\mathrm{kg\,m^2}$、$\omega_0=6.0\,\mathrm{rad/s}$，受恒阻力矩$M=-12\,\mathrm{N\,m}$，求减速到$2.0\,\mathrm{rad/s}$时转过的角度。
\workings
角加速度$\alpha=M/J=-4\,\mathrm{rad/s^2}$，
耗时$t=(2-6)/(-4)=1\,\mathrm s$，故
\[
\boxed{\theta=\omega_0t+\tfrac12\alpha t^2
=6-2=4\,\mathrm{rad}}.
\]
\crosscheck 也可由$\omega^2-\omega_0^2=2\alpha\theta$得
$(4-36)/(-8)=4$。

\sourceproblem{3.2.4}{69}{69}
\problemtext 两名各为$70\,\mathrm{kg}$的滑冰者，沿相距$10\,\mathrm m$的平行滑行路线以$6.5\,\mathrm{m/s}$相向经过；在交错瞬间抓住一条$10\,\mathrm m$绳的两端，随即绕绳中点转动。求各自对中心的角动量；若把绳缩短到$5\,\mathrm m$，求各自速率。
\workings
对称情形，初始每人距质心$r_0=5$ m，速度垂直半径，
\[
\boxed{L_0=mr_0v_0=70(5)(6.5)=2275\,\mathrm{kg\,m^2/s}}.
\]
忽略冰面摩擦，绳张力沿中心方向，不提供对中心的力矩。
最终$r_1=2.5$ m，且缩绳结束后径向速度为零，故
\[
\boxed{v_1=L_0/(mr_1)=13\,\mathrm{m/s}}.
\]
\crosscheck 距中心减半，速率加倍；缩绳者需做功，故动能不守恒。

\sourceproblem{3.2.5}{70}{70}
\problemtext 子弹沿水平方向射入下端系有质点球、可绕顶端水平轴转动的轻杆/球系统并滞留其中；判断射入的极短时间内与其后系统升高时各有什么守恒量。
\workings
极短碰撞中，轴对转轴的力矩为零，重力矩冲量可忽略，
所以\textbf{关于固定转轴的角动量守恒}：
\[
L_{\rm before}=L_{\rm after}.
\]
碰撞为完全非弹性，\emph{机械能不守恒}。
碰后上摆若忽略轴摩擦和空气阻力，并把地球引力势能计入，
\[
\boxed{K_{\rm after}+U_{\rm after}=K_{\rm top}+U_{\rm top}},
\]
即\textbf{机械能守恒}。
\crosscheck 碰撞阶段角动量守恒不推出转动动能守恒；提升阶段则需要将地球计入保守系统。

\sourceproblem{3.2.6}{70}{70}
\problemtext 质点系合外力恒为零时，判断合外力矩是否必为零，以及动量、角动量、机械能中哪一量必守恒。
\workings
合外力为零满足$\dot{\bm P}=0$，故
\[
\boxed{\bm P=\mathrm{const}}.
\]
但合外力矩可不为零：一对作用点不同、方向相反的外力构成力偶，
即$\sum\bm F=0$而$\sum\bm M\ne0$。故角动量不一定守恒；内摩擦可改变机械能。
\crosscheck 这里无额外“全部内力保守”条件，所以不能推出机械能守恒。

\sourceproblem{3.2.7}{70}{70}
\problemtext 花样滑冰运动员初始转动惯量$J_0$、角速度$\omega_0$；收臂后转动惯量降为原来的$1/3$。求最终角速度。
\workings
假设外力对竖直轴的力矩忽略，
$J_0\omega_0=(J_0/3)\omega_1$，故
\[
\boxed{\omega_1=3\omega_0}.
\]
\crosscheck 转动动能变为原来的三倍；收臂需人体肌肉做功，不能把动能也当作守恒量。

\sourceproblem{3.2.8}{70}{70}
\problemtext 飞轮以角速度$300\,\mathrm{rad/min}$旋转，转动惯量$J=5\,\mathrm{kg\,m^2}$，恒定阻力矩使其在$2$ s内停止。求阻力矩大小。
\workings
$\omega_0=300/60=5\,\mathrm{rad/s}$，
$\alpha=(0-5)/2=-2.5\,\mathrm{rad/s^2}$，故
\[
\boxed{|M|=J|\alpha|=12.5\,\mathrm{N\,m}}.
\]
\crosscheck 原题单位为$\mathrm{rad/min}$而不是每分钟转数；若误当$300\,\mathrm{r/min}$将引入额外$2\pi$因子。

\sourceproblem{3.2.9}{70}{70}
\problemtext 质量$m$的质点位置$\bm r=a\cos(\omega t)\bm i+
b\sin(\omega t)\bm j$，其中$a,b,\omega$均为正常量。求对坐标原点的合外力矩。
\workings
微分两次得$\bm a=-\omega^2\bm r$，故
$\bm F=-m\omega^2\bm r$为指向原点的中心力。
\[
\boxed{\bm M_O=\bm r\times\bm F=\bm0}.
\]
\crosscheck 即使轨迹为椭圆、速率不恒定，中心力对其中心的力矩仍为零。

\sourceproblem{3.2.10}{70}{70}
\problemtext 质量$m$、长度$L$的均匀细杆，可绕其\emph{中点}$O$处垂直于竖直面的水平定轴转动，两端分别固定$m$和$2m$的质点球。系统从与水平成$60^\circ$的静止位置释放。问杆转到水平时重力对系统关于$O$的合力矩大小。
\workings
杆的质心在$O$，故杆重力矩为零。两个球关于$O$的力臂均为$L/2$，但方向相反，合力矩
\[
\boxed{|M_O|=2mg\frac L2-mg\frac L2=\frac12mgL}.
\]
\crosscheck 此结论依“过中点$O$”与原图两端各一个质点球；若轴位置变化，力矩亦变化。

\sourceproblem{3.2.11}{70}{70}
\problemtext 刚体由$\omega_0=8\,\mathrm{rad/s}$在大小$8\,\mathrm{N\,m}$的恒定阻力矩作用下，$2$ s后降至$4\,\mathrm{rad/s}$。求角加速度及转动惯量。
\workings
\[
\boxed{\alpha=\frac{4-8}{2}=-2\,\mathrm{rad/s^2}},
\qquad
\boxed{J=\frac{|M|}{|\alpha|}=4\,\mathrm{kg\,m^2}}.
\]
\crosscheck 阻力矩为负，角加速度也为负，$J$必须为正。

\sourceproblem{3.2.12}{70}{70}
\problemtext 人站在转台中心，两臂侧平举时系统以$2\pi\,\mathrm{rad/s}$旋转，转动惯量为$6\,\mathrm{kg\,m^2}$；收臂后转动惯量为$2\,\mathrm{kg\,m^2}$。求末初转动动能之比。
\workings
关于竖直轴角动量守恒，$\omega_1=(6/2)\omega_0=3\omega_0$，
\[
\boxed{\frac{K_1}{K_0}
=\frac{\tfrac12J_1\omega_1^2}{\tfrac12J_0\omega_0^2}
=\frac{J_0}{J_1}=3}.
\]
\crosscheck 转动能增加三倍，由人体内部做功提供，而非外力矩。
% --- END textbook-ch03-blanks.tex ---

% --- BEGIN textbook-ch03-solutions-001-008.tex SHA ba140aadc0011f9fabc0aaf7ad3544ec6b8a292f ---
% 第三批：《大学物理简明教程》第5版，解答题3.3.1--3.3.8。
% 源书/PDF第70--72页。公式独立复算，原书图题模型单列，待二校。
\subsection{解答题：3.3.1--3.3.8}

\sourceproblem{3.3.1}{70}{70}
\problemtext 半径$R=0.30\,\mathrm m$的飞轮由静止开始做匀角加速度$\alpha=0.5\,\mathrm{rad/s^2}$的转动。求轮缘一点转过$240^\circ$时的切向加速度和法向加速度。
\workings
$\alpha$恒定，故角加速度仍为$0.5\,\mathrm{rad/s^2}$。
$\theta=240^\circ=4\pi/3$，运动学关系$\omega^2=2\alpha\theta$给
\[
\boxed{\omega=\sqrt{\frac{4\pi}{3}}\approx2.047\,\mathrm{rad/s}},
\qquad \boxed{a_\tau=R\alpha=0.15\,\mathrm{m/s^2}}.
\]
\[\boxed{a_n=R\omega^2=0.4\pi\approx1.257\,\mathrm{m/s^2}}.\]
\crosscheck 两个加速度分量分别为$R\alpha$与$R\omega^2$，不可混为一谈。

\sourceproblem{3.3.2}{70--71}{70--71}
\problemtext 半径$R=10\,\mathrm{cm}$、内壁光滑的半球形碗绕竖直对称轴匀速转动。一个小球相对碗静止，位置高于碗底$h=4\,\mathrm{cm}$。求所需角速度。
\workings
设小球到轴的水平距离为$r$、到球心的竖直投影绝对值$R-h$。
碗壁支持力$N$沿指向球心的法线，竖直方向平衡、水平向心：
\[
N\frac{R-h}{R}=mg,\qquad N\frac rR=m\omega^2r.
\]
排除$r=0$的底端特例，两式相除得
\[
\boxed{\omega=\sqrt{\frac{g}{R-h}}}
=\sqrt{\frac{9.8}{0.10-0.04}}
\approx12.78\,\mathrm{rad/s}.
\]
\crosscheck 位置几何$r^2=2Rh-h^2=(0.08\,\mathrm m)^2>0$；
当$h\to R^-$时所需角速度发散。

\sourceproblem{3.3.3}{70--71}{70--71}
\problemtext 均匀圆盘质量$m=1.1\,\mathrm{kg}$、转动惯量$J=\frac12mr^2$，
轮缘绳悬挂$m_1=1.0\,\mathrm{kg}$重物。施加外力矩使重物以$v_0=0.6\,\mathrm{m/s}$匀速上升，突然撤去外力矩，求重物开始反向运动前所经时间。
\workings
撤去外力矩后，重力使圆盘角速度减小。以重物向上为正、受力$T-m_1g=m_1a$，轮轴上由绳张力引起的减速力矩满足$-Tr=J(a/r)$。
联立或由整体有效质量得到
\[
a=-\frac{m_1g}{m_1+J/r^2}
=-\frac{m_1g}{m_1+m/2}.
\]
至速度为零所需时间
\[
\boxed{t=\frac{v_0}{|a|}
=\frac{v_0(m_1+m/2)}{m_1g}
=\frac{0.6(1+1.1/2)}{9.8}\approx0.0949\,\mathrm s}.
\]
\crosscheck $r$从有效转动惯量中消去，因圆盘惯量正比$r^2$；撤矩后的初速度仍向上，随后由重力制动。

\sourceproblem{3.3.4}{71}{71}
\problemtext 半径$R$、转动惯量$J$的水平光滑转台，以角速度$\omega_0$转动。质量$m$的人最初站在轴心，随后沿半径走到边缘，求新的角速度。
\workings
对竖直轴外力矩忽略，角动量守恒。
起初人的质点转动惯量可忽略，后为$mR^2$，因此
\[
J\omega_0=(J+mR^2)\omega_f,
\qquad \boxed{\omega_f=\frac{J}{J+mR^2}\omega_0}.
\]
\crosscheck $m\to0$恢复$\omega_0$；人移到更远处使系统角速率降低；非保守内功使动能未必守恒。

\sourceproblem{3.3.5}{71}{71}
\problemtext 质量$m=60\,\mathrm{kg}$、半径$R=0.25\,\mathrm m$的均匀圆盘飞轮以$900\,\mathrm{r/min}$旋转。杠杆左端为支点，支点到闸瓦接触点为$0.50\,\mathrm m$，闸瓦到施力端为$0.75\,\mathrm m$；闸瓦与轮缘动摩擦系数$\mu=0.4$。$(1)$若向下施力$F=100\,\mathrm N$，求停转时间与圈数；$(2)$若要求$2\,\mathrm s$内转速减半，求$F$。
\workings
飞轮转动惯量$J=\frac12mR^2=1.875\,\mathrm{kg\,m^2}$，
初角速度$\omega_0=900(2\pi)/60=30\pi\,\mathrm{rad/s}$。
闸杆视为无质量杠杆，关于左端支点力矩平衡
\[
N(0.50)=F(0.50+0.75),\qquad N=2.5F.
\]
摩擦制动力矩大小
$\tau=\mu NR=\mu(2.5F)R=0.25F\,\mathrm{N\,m}$
（$F$以牛顿的数值计）。
$(1)$ $F=100\,\mathrm N$时，
\[
|\alpha|=\tau/J=25/1.875=\frac{40}{3}\,\mathrm{rad/s^2},
\quad
\boxed{t_{\rm stop}=\frac{\omega_0}{|\alpha|}
\approx7.069\,\mathrm s}.
\]
停转角位移$\theta=\omega_0^2/(2|\alpha|)$，故
\[
\boxed{n=\frac{\theta}{2\pi}
=\frac{\omega_0^2}{4\pi|\alpha|}
\approx53.01\text{圈}}.
\]
$(2)$ $2$秒降到$\omega_0/2$要求
$|\alpha|=\omega_0/(4\,\mathrm s)$，于是
\[
\boxed{F=\frac{J(\omega_0/2)}{(2\,\mathrm s)(0.25\,\mathrm m)}
=\frac{1.875(30\pi)}{1.0}
\approx176.7\,\mathrm N}.
\]
\crosscheck 第二问的简化公式直接写为
$F=J[\omega_0/(2\Delta t)]/(\mu\,2.5R)$；
上式数值按$\Delta t=2\,\mathrm s$代入，见此式。
代数式中将$0.25$视为单位为米的力矩与外力比例系数，即$\tau/F=0.25\,\mathrm m$；制动时间与外力成反比。
\sourceproblem{3.3.6}{71}{71}
\problemtext 两个同轴均匀圆柱体固连，可绕水平轴无摩擦转动。大、小圆柱分别质量$M=10\,\mathrm{kg},m=4\,\mathrm{kg}$，半径$R=0.20\,\mathrm m,r=0.10\,\mathrm m$；绳分别在两柱绕向相反方向，吊质量$m_1=m_2=2\,\mathrm{kg}$的重物。求角加速度和两绳张力。
\workings
取$m_1$向下为正，$m_2$向上为正。合转动惯量
\[
J=\tfrac12MR^2+\tfrac12mr^2=0.22\,\mathrm{kg\,m^2}.
\]
加速度大小$a_1=\alpha R,\ a_2=\alpha r$，牛顿方程
\[
m_1g-T_1=m_1R\alpha,\quad
T_2-m_2g=m_2r\alpha,\quad
T_1R-T_2r=J\alpha.
\]
消去张力得
\[
\boxed{\alpha=\frac{m_1gR-m_2gr}
{J+m_1R^2+m_2r^2}
=\frac{2(9.8)(0.2-0.1)}{0.22+0.08+0.02}
=6.125\,\mathrm{rad/s^2}}.
\]
从而
\[
\boxed{T_1=m_1(g-R\alpha)=17.15\,\mathrm N},\qquad
\boxed{T_2=m_2(g+r\alpha)=20.825\,\mathrm N}.
\]
\crosscheck 直接代入$T_1R-T_2r=0.22\alpha$两边均为$1.3475\,\mathrm{N\,m}$。

\sourceproblem{3.3.7}{71--72}{71--72}
\problemtext 两个半径均为$R$、质量均为$m$的均匀圆盘状定滑轮，轻绳绕过两轮，两端悬挂质量分别为$m$和$2m$的重物。绳轮不滑动，求两滑轮间绳段的张力。
\workings
取$2m$侧向下为正。两轮转动惯量各为$J=\frac12mR^2$，整套系统有效惯量质量
$m+2m+2(J/R^2)=4m$。
\[
a=\frac{(2m-m)g}{4m}=\frac g4.
\]
设左侧张力$T_L=m(g+a)=\frac54mg$，右侧张力
$T_R=2m(g-a)=\frac32mg$。
左轮由中间张力$T_C$与左侧张力驱动：
$(T_C-T_L)R=Ja/R$，因此
\[
\boxed{T_C=T_L+\frac{Ja}{R^2}
=\frac54mg+\frac18mg=\frac{11}{8}mg}.
\]
\crosscheck 由右轮$(T_R-T_C)R=Ja/R$也得
$T_C=\frac32mg-\frac18mg=\frac{11}{8}mg$。

\sourceproblem{3.3.8}{71--72}{71--72}
\problemtext 长度$l$、质量$m$的均匀细杆在竖直平面内可绕中点$O$自由转动，初始水平且静止。一只质量同为$m$的小虫以垂直于细杆的速率$v_0$落在距$O$为$l/4$处，并沿杆向端点$A$爬行。求使杆角速度保持常值时，小虫需要的沿杆爬行速率。
\workings
\textbf{必须区分“恒定角速度”和“恒定爬行速率”，原题只要求前者。}
初始碰撞过程中对$O$轴角动量守恒：
\[
m v_0\frac l4
=\left(\frac{ml^2}{12}+m\frac{l^2}{16}\right)\omega_*,
\qquad \boxed{\omega_*=\frac{12v_0}{7l}}.
\]
碰后取由水平向下转过的角度$\theta=\omega_* t$为正，小虫距轴$r(t)$沿正方向向外，杆的质心位于轴上而重力矩为零。对整体应用角动量定理：
\[
\frac{d}{dt}\left[\left(\frac{ml^2}{12}+mr^2\right)\omega_*\right]
=mg r\cos\theta.
\]
因$\dot\omega_*=0$，只要$r>0$，可约去$mr$，得
\[
\boxed{v_{\rm crawl}(t)=\dot r
=\frac{g}{2\omega_*}\cos(\omega_*t)}.
\]
这给出\emph{保持角速度恒定所需的变速爬行规律}。
积分得
$r(t)=l/4+\frac{g}{2\omega_*^2}\sin(\omega_*t)$。
\crosscheck 在$0\le\theta\le\pi/2$时$\dot r\ge0$，可向$A$端爬行；
越过$\pi/2$后要求$\dot r<0$，故不能无限时间保持“持续向外爬行”和“恒定角速度”两项条件同时成立。若要求持续向外爬至端点，在$0\le\theta\le\pi/2$内须$l/4+g/(2\omega_*^2)\ge l/2$，即$\omega_*^2\le2g/l$，等价于$v_0^2\le49gl/72$（$v_0>0$）。到达端点即停止采用该爬行模型。
% --- END textbook-ch03-solutions-001-008.tex ---

% --- BEGIN textbook-ch03-solutions-009-015.tex SHA f839ef02d5725b283017215e93642cc39b4432f5 ---
% 第三批：《大学物理简明教程》第5版，解答题3.3.9--3.3.15。
% 来源：书内及用户PDF第72--73页；待独立学术二校。
\subsection{解答题：3.3.9--3.3.15}

\sourceproblem{3.3.9}{72}{72}
\problemtext 均匀细杆质量$m_1$、长度$l$，上端绕水平轴$O$，初始竖直静止。质量$m_2$的小球水平以速率$v_0$击中距轴$d=\frac34l$处，碰后黏附在杆上。求碰撞刚结束时的角速度。
\workings
短时碰撞中，支点冲量对$O$无力矩，重力冲量矩忽略，故
\[
m_2v_0d=(J_{\rm rod}+m_2d^2)\omega,\qquad
J_{\rm rod}=\frac13m_1l^2.
\]
代入$d=3l/4$得到
\[
\boxed{\omega=
\frac{\frac34m_2lv_0}
{\frac13m_1l^2+\frac9{16}m_2l^2}
=\frac{36m_2v_0}{l(16m_1+27m_2)}}.
\]
\crosscheck $m_2\to0$时$\omega\to0$；该非弹性碰撞不能用机械能守恒计算瞬间角速度。

\sourceproblem{3.3.10}{72}{72}
\problemtext 光滑桌面上的$m_1=5\,\mathrm{kg}$物块通过轻绳连接垂挂的$m_2=20\,\mathrm{kg}$物块，绳跨过半径$r=0.10\,\mathrm m$、质量$M=1.5\,\mathrm{kg}$的均匀圆盘定滑轮，不打滑。求两物体共同的加速度。
\workings
设两侧绳张力为$T_1,T_2$，沿运动方向有
$T_1=m_1a,\ m_2g-T_2=m_2a$。
轮转动方程
\[
(T_2-T_1)r=J\frac ar,\qquad J=\tfrac12Mr^2.
\]
相加得
\[
m_2g=\left(m_1+m_2+\frac J{r^2}\right)a,
\quad\boxed{
a=\frac{m_2g}{m_1+m_2+M/2}
=\frac{20(9.8)}{5+20+0.75}
\approx7.612\,\mathrm{m/s^2}}.
\]
\crosscheck $M\to0$恢复理想无惯量滑轮加速度$20g/25$；$r$因$J\propto r^2$消去。

\sourceproblem{3.3.11}{72}{72}
\problemtext 一根长$l$、质量$m$的均匀细杆绕左端水平光滑轴$O$在竖直面内自由转动；杆从\emph{水平位置}静止释放。求初始角加速度及杆从水平位置向下转过$\theta$时的角速度。
\workings
取向下转动为正。初始重力矩为$mg(l/2)$，杆对端轴惯量$J=ml^2/3$，所以
\[
\boxed{\alpha_0=\frac{mg(l/2)}{ml^2/3}=\frac{3g}{2l}}.
\]
转角$\theta$后杆的质心下降$(l/2)\sin\theta$，重力势能损失转为转动动能：
\[
mg\frac l2\sin\theta=\frac12\frac{ml^2}{3}\omega^2.
\]
故
\[
\boxed{\omega(\theta)=\sqrt{\frac{3g\sin\theta}{l}}}
\quad(0\le\theta\le\pi\text{，首次下摆到另一侧水平位置}).
\]
\crosscheck $\theta=0$时$\omega=0$；当$\theta=\pi/2$速率最大，且$\alpha$此时为零。

\sourceproblem{3.3.12}{72--73}{72--73}
\problemtext 质量$m$的质点悬挂于天花板固定点$B$，绳长$l$，以水平圆轨道半径$r<l$、角速度$\omega$做匀速圆周运动，水平圆心为$O$。求对$O$与$B$的角动量、重力及张力对两点的力矩，并判断角动量守恒。
\workings
令圆轨道位于$z=0$平面，$B$在$O$上方
$h=\sqrt{l^2-r^2}$。取质点瞬时位置$\bm r_O=r\bm e_r$、速度$\bm v=\omega r\bm e_\varphi$。
则
\[
\boxed{\bm L_O=mr^2\omega\bm e_z},\qquad
\bm r_B=r\bm e_r-h\bm e_z,
\quad
\boxed{\bm L_B=m\omega r(h\bm e_r+r\bm e_z)}.
\]
因此$|\bm L_O|=mr^2\omega$、$|\bm L_B|=m\omega rl$。
张力$T$指向$B$，竖直平衡$Th/l=mg$。相对$O$，重力力矩
\[
\bm M_{g,O}=mgr\,\bm e_\varphi,
\]
张力关于$O$的力矩
$\bm M_{T,O}=-Tr(h/l)\bm e_\varphi=-mgr\,\bm e_\varphi$，两者相互抵消。
故$\boxed{\dot{\bm L}_O=0}$。
相对$B$，张力力矩为零，而重力力矩
$\boxed{\bm M_{g,B}=mgr\,\bm e_\varphi}$；因$\bm e_r$随时间转动，
$\dot{\bm L}_B=m\omega^2rh\,\bm e_\varphi$，利用$\omega^2h=g$与力矩一致，故
$\bm L_B$一般\emph{不守恒}。
\crosscheck 注意$\bm L_B$的\emph{模}虽恒定，\emph{矢量}的方向在变；“模不变”不能代替“角动量守恒”。

\sourceproblem{3.3.13}{72--73}{72--73}
\problemtext 质量$M$、半径$R$的均匀飞轮以$\omega$转动，一块质量$m$、相对飞轮可视为边缘点质量的小碎片瞬时脱离，离开时速度恰竖直向上。求碎片最高上升高度与余下飞轮刚脱离后的角速度、角动量和转动动能。
\workings
碎片初始速率$v=\omega R$，忽略空气阻力，
\[
\boxed{h_{\max}=\frac{\omega^2R^2}{2g}}.
\]
假定碎裂过程没有施加额外的切向冲量矩，剩余各材料点速度瞬时连续，因此飞轮余部刚脱离后仍有
$\boxed{\omega_{\rm rem}=\omega}$。
以碎片为位于轮缘的点质量近似，原惯量$J_0=\frac12MR^2$，
余部惯量$J_{\rm rem}=J_0-mR^2$，从而
\[
\boxed{L_{\rm rem}=\left(\tfrac12M-m\right)R^2\omega},
\qquad
\boxed{K_{\rm rem}=\frac12\left(\tfrac12M-m\right)R^2\omega^2}.
\]
\crosscheck 上式是“飞轮剩余部分”的量，而非加上已脱离碎片的总角动量；点碎片近似要求$m<M/2$且碎片脱离不改变余部其他质量点瞬时速度。

\sourceproblem{3.3.14}{73}{73}
\problemtext 质量$m$、半径$R$、质量集中于轮缘的车轮起初静止，可绕轴$O$自由转动。质量$m_0$的子弹以水平速率$v_0$射入轮缘并留在轮中，入射高度与入射处径向满足图中的角$\theta$。求碰后角速度，并求最终系统转动动能占子弹初动能的比例。
\workings
初始车轮静止，子弹对轴的角动量大小
$L_i=m_0v_0R|\sin\theta|$，轮缘惯量$J_w=mR^2$；
子弹射入后系统惯量$J_f=(m+m_0)R^2$。
\[
\boxed{|\omega|=\frac{m_0v_0|\sin\theta|}{(m+m_0)R}}.
\]
相应最终转动动能
\[
K_f=\frac12J_f\omega^2
=\frac{m_0^2v_0^2\sin^2\theta}{2(m+m_0)},
\]
因此
\[
\boxed{\frac{K_f}{K_i}=\frac{m_0}{m+m_0}\sin^2\theta},
\qquad K_i=\frac12m_0v_0^2.
\]
\crosscheck 完全非弹性碰撞中该比值不大于$1$；$\theta=0$（沿径向射入）时对子弹角动量的贡献为零。

\sourceproblem{3.3.15}{73}{73}
\problemtext 弹簧$k=2.0\,\mathrm{N/m}$经一均匀定滑轮与质量$M=6.0\,\mathrm{kg}$的悬挂物体连接；滑轮转动惯量$J=0.5\,\mathrm{kg\,m^2}$、半径$R=0.30\,\mathrm m$，开始弹簧未伸长且物体静止。物体下落$h=0.40\,\mathrm m$时求其速率，忽略摩擦和绳滑动。
\workings
下落时弹簧同样伸长$h$，绳速率$v=R\omega$。
以物体、滑轮与弹簧为系统，机械能守恒：
\[
Mgh=\tfrac12Mv^2+\tfrac12J\frac{v^2}{R^2}+\tfrac12kh^2.
\]
整理得到
\[
\boxed{v=\sqrt{\frac{2Mgh-kh^2}{M+J/R^2}}}
=\sqrt{\frac{2(6)(9.8)(0.4)-2(0.4)^2}{6+0.5/(0.3)^2}}
\approx2.011\,\mathrm{m/s}.
\]
\crosscheck 在理想模型中，$0<h<2Mg/k$时速率为正，$h=2Mg/k$为可达转折点；转动惯量增大使速率降低，量纲为$\mathrm{m/s}$。
% --- END textbook-ch03-solutions-009-015.tex ---
